% Separate post-CP32 candidate; B91 dependency remains separately audited.
\subsection{Actual diagonal fields remove one residual-wild crossing}

Throughout this subsection \(K\) is complete discretely valued of
mixed characteristic \((0,5)\), with algebraically closed residue
field, \(e_K=20\), and \(\mu_5\subset K\). The common tame good
field \(L/K\) has degree 24, so \(v_L(5)=480\). All point fields,
quotients and matrix images below arise from the stated actual
torsion comparisons. Breaks use \(b(g)=v(g\varpi-\varpi)-1\).

\begin{lemma}[The exchanged \(25\)-comparison already has full first image]
\label{lem:B-RD92-tame25}
\status{Candidate; complete proof and exact supplement pending independent audits.}
Suppose two whole potentially supersingular height-two groups have
isomorphic actual \(25\)-torsion, rational first-level torsion over
\(L\), parameters \(25,75\), and exchanged high lines
\(\mathbf F_5e_1,\mathbf F_5e_2\). Assume their normalized residual
tame exponents are \((0,6)\), or \((12,18)\), in this basis.
Their common \(25\)-division image over \(L\) is
\(1+5M_2(\mathbf F_5)\). Its four coordinate cyclic characters
have breaks
\[
 \begin{array}{c|rrrr}
 \text{character}&A_{21}&A_{22}&A_{12}&A_{11}\\\hline
 \text{break}&6&24&66&72.
 \end{array}                                           \tag{B-RD92.1}
\]
No actual \(125\)-comparison is assumed in this lemma.
\end{lemma}
\begin{proof}
Let \(V\subset M_2(\mathbf F_5)\) be the actual additive image.
Choose a second-level point above \(e_i\), with field \(E_i\).
First-level rationality makes \(E_i/L\) Galois with translation
group \(Ve_i\). Its degree is at most 25. The inverse Newton
polygons in both models are single segments: both parameters exceed
20, including the high-target inverse. The coordinate valuations
over \(L\), in the two models, are
\[
  (19/5,3/5)\quad(i=1),\qquad
  (1/5,9/5)\quad(i=2).                                \tag{B-RD92.2}
\]
Thus each degree is at least five. If it were five, all corresponding
normalized coordinate valuations would be prime to five, and the
exact B60 break formula would apply. There are three possible types
of its translation line: the first high line, the second high line,
or a line different from both. For \(i=1\), the two exact breaks
would respectively be
\[
          (456,72),\qquad(6,222),\qquad(6,72);
                                                        \tag{B-RD92.3}
\]
for \(i=2\), they would be
\[
          (474,66),\qquad(24,216),\qquad(24,66).
                                                        \tag{B-RD92.4}
\]
For example the translating high/low valuations over a degree-five
field are \((475,25)\) in the first model and \((225,75)\) in the
second. Subtracting the two coordinate values in (B-RD92.2), scaled
by five, gives these rows. None agrees. Hence both evaluations
\(V\to\mathbf F_5^2\) are surjective.

Tame conjugation has distinct weights 0, 6 and 18 on the diagonal
plane and the two off-diagonal axes. Surjectivity of both evaluations
therefore puts \(E_{12},E_{21}\) in \(V\), and makes both diagonal
projections nonzero. Consequently \(\dim V\ge3\). If equality
held, the diagonal part would be \(k_5\operatorname{diag}(a,b)\),
with \(ab\ne0\); the two degree-25 point fields would share its
actual diagonal cyclic quotient. Here \(k_5=\mathbf F_5\).

The complete shared-diagonal proof of Proposition B87 then applies
with no higher torsion level: its four coordinate rows are
\[
\begin{array}{c|c|r|r|r|c}
 \text{field}&\text{model}&A&D_0&D_1&H\\\hline
 E_1&1&95&125&2375&\text{diagonal}\\
 E_1&2&15&375&1125&\text{off-diagonal}\\
 E_2&1&5&125&2375&\text{off-diagonal}\\
 E_2&2&45&375&1125&\text{diagonal}.
\end{array}                                             \tag{B-RD92.5}
\]
All four displacements are below \(v_{E_i}(5)=12000\), and every
\(A\) has five-adic order one. The diagonal cyclic break is a
positive multiple of 24, while the off-diagonal breaks are 6 and
18 modulo 24. In \(E_1\), an off-diagonal last line would give its
larger cyclic break 6, impossible; hence its diagonal break is 72.
In \(E_2\), a diagonal last line gives diagonal break 24, whereas
an off-diagonal last line gives off-diagonal break 66, which could
not exceed the shared diagonal break 72. This excludes dimension
three, using exactly the fixed-coordinate lemma B87 and the same
actual quotient field. Thus \(V=M_2(\mathbf F_5)\).

Finally the complete point-field calculation in B88 uses only this
actual full \(25\)-image and (B-RD92.5). To spell out its finite
alternatives, in \(E_1\) the last line is diagonal, with cyclic
break 72; the smaller break is \(t\in\{6,30,54\}\), and the
larger lower break is \(360-4t\). At \(t=30\) the compulsory
prime exponent is 135; at \(t=54\) the other row would require
an exponent \(71<95\). Hence only \((t,u)=(6,336)\) remains.
In \(E_2\), a diagonal last line gives \((18,48)\), whose forced
prime term changes at 165 instead of 1125. An off-diagonal last
line gives cyclic break 66 and \(t\in\{24,48\}\). The latter
case forces a change at 215 instead of 2375. Thus
\((t,u)=(24,234)\). The corresponding two pairs of cyclic
breaks are \((A_{11},A_{21})=(72,6)\) and
\((A_{22},A_{12})=(24,66)\), proving (B-RD92.1).
\end{proof}

\begin{lemma}[The actual image in either residual-wild crossing]
\label{lem:B-RD92-parahoric}
\status{Candidate; complete proof and exact supplement pending independent audits.}
Suppose an actual \(125\)-comparison is in one of the B91
crossing rows \((5,45)\) or \((15,95)\), with common high line
\(H=\mathbf F_5e_1\). Its full wild image is
\[
 W=\left\{
 \begin{pmatrix}1+5a&b\\25c&1+5d\end{pmatrix}:
 a,d\in\mathbf Z/25,\ b\in\mathbf Z/125,
 c\in\mathbf F_5\right\},\qquad |W|=5^8.            \tag{B-RD92.6}
\]
In particular its first and terminal layers are respectively
the upper-triangular three-space and all of \(M_2(\mathbf F_5)\).
Let \(R/L\) be the original first-level division field, and let
\(D_{22}/L\) be the cyclic field defined by
\((\rho_{22}-1)/5\bmod5\). Their breaks are respectively
\[
 \begin{array}{c|cc}
 (B_1,B_2)&b(R/L)&b(D_{22}/L)\\\hline
 (5,45)&66&24\\
 (15,95)&6&72.
 \end{array}                                           \tag{B-RD92.7}
\]
The two fields in each row are distinct actual subfields.
\end{lemma}
\begin{proof}
Quotient both groups by their common high line. B91 gives an
actual \(25\)-comparison with rational first-level torsion and
exchanged parameters \(25,75\). The preceding lemma applies.
In the original integral basis the quotient lattice has basis
\((e_1/5,e_2)\), so its matrices are
\[
       \begin{pmatrix}\rho_{11}&5\rho_{12}\\
                   \rho_{21}/5&\rho_{22}\end{pmatrix}.
                                                        \tag{B-RD92.8}
\]
Rationality of the quotient first-level module therefore forces
\(25\mid\rho_{21}\) on the entire original wild image. It lies
in the displayed group (B-RD92.6).

The full quotient \(25\)-image makes the four additive characters
\[
 \rho_{12}\bmod5,\quad
 (\rho_{11}-1)/5\bmod5,\quad
 (\rho_{22}-1)/5\bmod5,\quad
 \rho_{21}/25\bmod5                                  \tag{B-RD92.9}
\]
jointly surjective. They are additive directly from the matrix
product and \(25\mid\rho_{21}\). Choose an integral eigenbasis
for an actual order-24 tame lift. The first congruence layer is
upper triangular, and its diagonal projection is surjective:
set the first character of (B-RD92.9) to zero while prescribing
the next two. Spectral projection gives the full diagonal plane.
Also, an element with nonzero residual upper entry has fifth power
\(1+5uE_{12}\pmod{25}\). This follows either by multiplying
the triangular matrices or by summing the five upper entries;
the terms from their diagonal corrections have coefficient ten
and vanish modulo 25. Thus the first layer is the full upper
triangular three-space.

If the last character in (B-RD92.9) vanished on the terminal layer,
its restriction to the first congruence subgroup would factor
through that upper-triangular first layer. Its tame weight is the
lower off-diagonal weight, distinct from both the diagonal and
upper weights, so that restriction would vanish. It would then
factor through the original residual upper character, again of
the other weight, and would be zero altogether. This contradicts
the joint surjectivity. Thus the terminal layer has a nonzero lower
entry. Tame projection supplies \(E_{21}\), and fifth powers
already put the whole first layer in the terminal layer. The latter
is therefore \(M_2(\mathbf F_5)\). The total order is
\(5^{1+3+4}=5^8\), the order of the containing group in
(B-RD92.6); equality follows.

It remains to identify the actual cyclic fields, including the
index order. In the \((5,45)\) row, the quotient with parameter25
is the first model and its high line is the first coordinate of
(B-RD92.8). Thus the original residual upper character is its
\(A_{12}\), and the original lower diagonal is its \(A_{22}\).
Their breaks in (B-RD92.1) are 66 and 24. In the \((15,95)\)
row, the parameter25 quotient is the second model: its high line
is the old first-level image, the second coordinate of
(B-RD92.8). Reordering the quotient basis to put this high line
first sends the original upper residual character to \(A_{21}\)
and the original lower diagonal to \(A_{11}\). Their breaks are
6 and 72. These are distinct independent characters on the actual
full additive quotient image, proving (B-RD92.7) and distinctness.
\end{proof}

\begin{lemma}[Two coordinates with the same high-translation line]
\label{lem:B-RD92-common-high}
\status{Candidate; complete proof and exact supplement pending independent audits.}
Let \(R/L\) and \(E/R\) be actual extensions in the crossing
setting, with \([R:L]=5\), \(\operatorname{Gal}(E/R)\simeq
\mathbf F_5^2\), and let \(b_1<b_2\) be prime to five,
\(1\le b_i\le19\). Suppose two formal coordinates in \(E\)
have the same high-displacement line and data
\[
 A_i=5b_i,\qquad D_{0,i}=125b_i,\qquad
 D_{1,i}=15000-625b_i,\qquad v_E(5)=60000.             \tag{B-RD92.10}
\]
Then the last ramification line is the common high line, the two
lower breaks are
\[
 t=24b_1,\qquad u=15000-630b_2-96b_1,                 \tag{B-RD92.11}
\]
and the cyclic quotient killing the high line has break \(24b_1\).
\end{lemma}
\begin{proof}
All displacement values are strictly below the mixed-characteristic
cutoff. Since \(v_5(A_i)=1\), the same fixed-coordinate expansion
as in B87 has first possible change at
\(\min(A_i+5r,J_i+r)\), where \(J_i>A_i\) is the first nonzero
prime-to-five exponent, possibly infinite. A strict smaller value
can cancel only when \(J_i=A_i+4r\).

There cannot be a single ramification break \(r\). Its leading
coefficient map on \(\mathbf F_5^2\) has a two-dimensional image.
A first change strictly below \(D_{0,i}\) would be a nonzero
additive polynomial of degree at most five vanishing on all 25
leading coefficients, impossible. If the first change were
\(D_{0,i}\) with a unique monomial, it could not vanish exactly
on the high line. Thus a tie at \(D_{0,i}\) is required, giving
\(r=(D_{0,i}-A_i)/5=24b_i\) in both models, a contradiction.

Write the two breaks as \(t<u\), with last line \(T\). If
\(T\) differed from the common high line, the coefficient lemma
B87 would give the same larger cyclic quotient break \(24b_i\)
in both models. Hence \(T\) is the common high line. Put
\(r_i=24b_i\), so \(r_1<r_2\). At the smaller break, every
element outside the high line has displacement \(D_{0,i}\).
Consequently
\[
 \begin{array}{c|c|c}
 &J_i&\text{forced relation at }u\\\hline
 t<r_i&A_i+4t&u+4t=15000-630b_i\\
 t>r_i&D_{0,i}-t&u-t=15000-750b_i.
 \end{array}                                           \tag{B-RD92.12}
\]
In both rows the prime term is uniquely first at \(u\). If both
models use the same strict row, their parameters are equal. If
\(r_1<t<r_2\), subtracting the two relations gives
\(t=150b_1-126b_2<24b_1\), impossible.

If \(t=r_2\), the first row's model, namely model1 in the second
strict case, gives \(u=15000-750b_1+24b_2\). In model2 one must
have \(J_2\ge101b_2\). At \(u\), its leading divisible term
exceeds the required displacement by
\[
 (5b_2+5u)-D_{1,2}=750(80-5b_1+b_2)>0.
\]
It therefore requires \(J_2=D_{1,2}-u=750b_1-649b_2\), which
would force \(b_1\ge b_2\). Thus \(t=r_2\) is impossible.
The only remaining value is \(t=r_1\); model2 is then in the
first strict row of (B-RD92.12), giving (B-RD92.11). The exact
cyclic quotient killing the last line has the smaller break \(t\),
as follows directly from the quotient-average formula. This proves
the assertion without assuming that every formal compatible profile
is realized by a field or a group.
\end{proof}

\begin{theorem}[The \((15,95)\) crossing is impossible]
\label{thm:B-RD92-crossing}
\status{Candidate; complete proof and exact supplement pending independent audits.}
The residual-wild crossing row \((B_1,B_2)=(15,95)\) cannot occur
in an actual different-special-character \(125\)-comparison.
Consequently a remaining crossing must have ordered parameters
\((5,45)\), and then has the full actual image (B-RD92.6),
with residual field break66 and diagonal field break24 as above.
\end{theorem}
\begin{proof}
Take an actual primitive \(25\)-point \(P\) with \(5P=e_2\)
outside the original common high line, and put \(E=L(P)\).
The full group (B-RD92.6) modulo25 has order \(5^4\). Its orbit
on this point has size \(5^3\): the upper entry modulo25 and
the lower diagonal correction modulo five vary independently.
Thus \([E:L]=125\), and \(R\subset E\). Over \(R\), the
congruence image is the additive upper triangular three-space;
evaluation on \(e_2\) is surjective with kernel
\(\mathbf F_5E_{11}\). Therefore \(E/R\) is an actual Galois
extension of degree25 with every first-level translation present.
The low inverse coordinates have \(v_L(T(P))=b_i/25\), while
the translating points have values \(b_i\) and \(120-5b_i\).
Normalization by \([E:L]=125\) gives exactly (B-RD92.10).

In the translation coordinates for \(E/R\), the high line is
the upper entry and the quotient killing it is the lower diagonal
character. Hence its actual cyclic field is \(D_{22}R/R\).
In the \((15,95)\) row one has \((b_1,b_2)=(3,19)\), so
the preceding lemma gives
\[
                         b(D_{22}R/R)=72.             \tag{B-RD92.13}
\]
But (B-RD92.7) gives two distinct cyclic fields over \(L\) with
breaks \(b(R/L)=6\) and \(b(D_{22}/L)=72\). Their actual
elementary abelian compositum has upper breaks6 and72, and hence
lower breaks6 and \(6+5(72-6)=336\). Restriction to its subgroup
fixing \(R\), which is the last line, gives
\[
                         b(D_{22}R/R)=336,             \tag{B-RD92.14}
\]
a contradiction. The conversions are precisely subgroup lower
numbering and quotient upper numbering, as in the full norm and
Herbrand proofs already used in B88--B90. No relation between
unrelated abstract fields is used.

For the other row the same test gives24 over \(R\), consistently
with breaks66 and24 over \(L\). The theorem makes no exclusion
or existence claim for that row.
\end{proof}

\begin{remark}[The remaining actual crossing]
\label{rem:B-RD92-scope}
\status{Candidate scope; no unrestricted least-index closure.}
The lower row is excluded by an actual cyclic field comparison,
not by the size of a hypothetical matrix group. The row \((5,45)\)
remains open despite the full necessary parahoric image. The
both-large residual-wild alternatives of B91 and the irreducible
residual branch are unchanged. The unrestricted index at
\(p=5,e=20\) therefore remains in \(\{2,3\}\).

The proof explicitly reuses only those parts of B87 and B88 that
concern the actual \(25\)-division image, and supplies the missing
degree-five point-field exclusions needed to obtain that image
without a \(125\)-comparison of the high quotients. The old level
is not silently restored after quotienting. The retained standalone
program \path{code/reduction_descent/verify_residual_wild_crossing.py}
reconstructs all 12 degree-five translation lines, the four possible
three-dimensional diagonal lines, the five relevant point-field
coefficient cases and the separate excluded diagonal case. It checks
625 actual triangular fifth powers and 3,125 quotient matrix-digit
rows, then tests 2,138,160 necessary common-high break pairs. The
surviving necessary profiles are exactly \((24,9234)\) and
\((72,2742)\); the actual diagonal-field comparison excludes only
the second crossing. The first author run passed at
03:16:05.929610 UTC on 20 November 2026 in
1.205554521002341 seconds. The complete program, output and
pre-certificate source are retained; no executed computation failed.
These finite checks supplement the actual field and coefficient
proofs and imply no realization of a surviving numerical profile.
The complete internal B60, B69, B87, B88, B90 and B91 dependency
proofs were reread; no fresh external original-source opening is
claimed in this candidate.
\end{remark}
